\documentclass[12pt,a4paper]{article}
\usepackage[top=2.5cm,bottom=2.5cm,left=2.5cm,right=2.5cm,headheight=15pt]{geometry}
\usepackage{amsmath,amssymb}
\usepackage{tikz}
\usetikzlibrary{shapes,arrows,arrows.meta,positioning,calc}
\usepackage{booktabs,array}
\usepackage{xcolor}
\usepackage{enumitem}
\usepackage{fancyhdr}
\usepackage{titlesec}
\usepackage{mdframed}
\usepackage{tcolorbox}
\tcbuselibrary{skins,breakable}
\usepackage{hyperref}

%% ---- Colors ----
\definecolor{folblue}{RGB}{0,82,165}
\definecolor{lightblue}{RGB}{220,235,255}
\definecolor{darkorange}{RGB}{180,70,0}
\definecolor{lightorange}{RGB}{255,235,210}
\definecolor{darkgreen}{RGB}{20,110,20}
\definecolor{lightgreen}{RGB}{200,240,205}
\definecolor{darkred}{RGB}{150,20,20}
\definecolor{lightred}{RGB}{255,210,210}
\definecolor{lightgray}{RGB}{242,242,242}

%% ---- tcolorbox styles ----
\newtcolorbox{qbox}[2][]{enhanced,breakable,
  colback=lightorange,colframe=darkorange,arc=4pt,
  title={\textbf{#2}},fonttitle=\bfseries\color{white},
  attach boxed title to top left={yshift=-2mm,xshift=4mm},
  boxed title style={colback=darkorange,arc=3pt},#1}

\newtcolorbox{solbox}[2][]{enhanced,breakable,
  colback=lightblue,colframe=folblue,arc=4pt,
  title={\textbf{#2}},fonttitle=\bfseries\color{white},
  attach boxed title to top left={yshift=-2mm,xshift=4mm},
  boxed title style={colback=folblue,arc=3pt},#1}

\newtcolorbox{keybox}[2][]{enhanced,breakable,
  colback=lightgreen,colframe=darkgreen,arc=4pt,
  title={\textbf{#2}},fonttitle=\bfseries\color{white},
  attach boxed title to top left={yshift=-2mm,xshift=4mm},
  boxed title style={colback=darkgreen,arc=3pt},#1}

\newtcolorbox{algobox}[2][]{enhanced,breakable,
  colback=lightgray,colframe=gray!60,arc=4pt,
  title={\textbf{#2}},fonttitle=\bfseries,
  attach boxed title to top left={yshift=-2mm,xshift=4mm},
  boxed title style={colback=gray!60,arc=3pt},#1}

\newtcolorbox{warnbox}[2][]{enhanced,breakable,
  colback=lightred,colframe=darkred,arc=4pt,
  title={\textbf{#2}},fonttitle=\bfseries\color{white},
  attach boxed title to top left={yshift=-2mm,xshift=4mm},
  boxed title style={colback=darkred,arc=3pt},#1}

%% ---- Page style ----
\pagestyle{fancy}\fancyhf{}
\lhead{\color{folblue}\textbf{Quantum Attacks \& ECC --- CSE 717 InfoSec}}
\rhead{\color{folblue}Exam: 17 Jun 2026}
\cfoot{--- \thepage\ ---}
\renewcommand{\headrulewidth}{0.4pt}

%% ---- Section style ----
\titleformat{\section}{\large\bfseries\color{folblue}}{}{0em}{}[\titlerule]
\titleformat{\subsection}{\normalsize\bfseries\color{darkorange}}{}{0em}{}

\begin{document}

\begin{center}
  {\LARGE\bfseries\color{folblue} Quantum Attacks \& ECC}\\[4pt]
  {\large CSE 717 Information Security --- Past-Paper Solutions (2024)}\\[2pt]
  {\small Source: Lc\#6 (ECC), Lc\#7 (Quantum Attacks) \quad Yield: 1/5 years (2024 only), 5 marks}
\end{center}
\vspace{0.5em}

%% ================================================================
\section{Algorithm Cheat Sheet}
%% ================================================================

\begin{algobox}{ECC Curve Definition}
An elliptic curve over $\mathbb{F}_p$ is defined by:
\[
  y^2 \equiv x^3 + ax + b \pmod{p}
\]
Each valid point $(x,y)$ satisfies this equation mod $p$. The point at infinity $\mathcal{O}$ acts as the identity (0 in addition).
\end{algobox}

\vspace{0.5em}

\begin{algobox}{Point Doubling Formula --- $2P$ where $P = Q = (x_1, y_1)$}
\textbf{Step 1: Compute slope} $\lambda$
\[
  \lambda = \frac{3x_1^2 + a}{2y_1} \pmod{p}
\]
\textbf{Note:} Division means multiply by the modular inverse. Use $ax \equiv 1 \pmod{p}$ to find $a^{-1}$.

\medskip
\textbf{Step 2: Compute} $x_3$
\[
  x_3 = \lambda^2 - 2x_1 \pmod{p}
\]
\textbf{Step 3: Compute} $y_3$
\[
  y_3 = \lambda(x_1 - x_3) - y_1 \pmod{p}
\]
\textbf{Result:} $2P = (x_3, y_3)$

\medskip
\textbf{Negative mod trick:} If result is negative, add multiples of $p$ until positive.
$(-a) \bmod p = (p - (a \bmod p)) \bmod p$
\end{algobox}

\vspace{0.5em}

\begin{algobox}{Point Addition Formula --- $P + Q$ where $P \neq Q$, $P=(x_1,y_1)$, $Q=(x_2,y_2)$}
\[
  \lambda = \frac{y_2 - y_1}{x_2 - x_1} \pmod{p}
\]
Then $x_3 = \lambda^2 - x_1 - x_2 \pmod{p}$, \quad $y_3 = \lambda(x_1 - x_3) - y_1 \pmod{p}$
\end{algobox}

\vspace{0.5em}

\begin{keybox}{Quantum Attack Summary}
\begin{tabular}{@{}lll@{}}
\toprule
\textbf{Algorithm} & \textbf{Target} & \textbf{Impact} \\
\midrule
Shor's Algorithm & RSA (integer factorization) & \textbf{Completely breaks} RSA \\
Shor's Algorithm & ECC (discrete log, ECDLP) & \textbf{Completely breaks} ECC \\
Grover's Algorithm & Symmetric (AES, DES) & Halves key strength (quadratic speedup only) \\
\bottomrule
\end{tabular}

\medskip
\textbf{Key distinction:} Shor's gives \emph{polynomial} speedup (exponential improvement) for public-key.
Grover's gives only $\sqrt{\cdot}$ speedup (quadratic) for symmetric --- manageable by doubling key size.
AES-256 $\to$ still quantum-resistant. AES-128 $\to$ drops to $\sim$64-bit equivalent security.
\end{keybox}

\vspace{1em}

%% ================================================================
\section{2024 Exam (1/5 years --- 5 marks total)}
%% ================================================================

\subsection{Q6(a) --- 3 marks: ECC Point Doubling}

\begin{qbox}{2024 Q6(a) [3 marks]}
  Elliptic curve $E\colon y^2 = x^3 + 2x + 3$ over $\mathbb{F}_{19}$.
  Compute $2P$ for $P = (3, 6)$.
\end{qbox}

\begin{solbox}{Solution --- Point Doubling (Case B: $P = Q$)}
\textbf{Given:} $x_1 = 3$, $y_1 = 6$, $a = 2$, $p = 19$

\bigskip
\textbf{Step 1: Compute slope $\lambda$}
\[
  \lambda = \frac{3x_1^2 + a}{2y_1} \pmod{19}
  = \frac{3(3)^2 + 2}{2(6)} \pmod{19}
  = \frac{27 + 2}{12} \pmod{19}
  = \frac{29}{12} \pmod{19}
\]

Convert numerator: $29 \bmod 19 = 10$

Find $12^{-1} \bmod 19$: need $12k \equiv 1 \pmod{19}$
\[
  12 \times 8 = 96 = 5 \times 19 + 1 \implies 96 \equiv 1 \pmod{19}
\]
So $12^{-1} \equiv 8 \pmod{19}$.

\[
  \lambda = 10 \times 8 = 80 \pmod{19} = 80 - 4(19) = 80 - 76 = \boxed{4}
\]

\bigskip
\textbf{Step 2: Compute $x_3$}
\[
  x_3 = \lambda^2 - 2x_1 \pmod{19}
  = 4^2 - 2(3) \pmod{19}
  = 16 - 6 \pmod{19}
  = \boxed{10}
\]

\bigskip
\textbf{Step 3: Compute $y_3$}
\[
  y_3 = \lambda(x_1 - x_3) - y_1 \pmod{19}
  = 4(3 - 10) - 6 \pmod{19}
  = 4(-7) - 6 \pmod{19}
  = -28 - 6 \pmod{19}
  = -34 \pmod{19}
\]
Convert: $-34 + 2(19) = -34 + 38 = \boxed{4}$

\bigskip
\[
  \therefore \quad \mathbf{2P = (10,\ 4)}
\]

\textit{Verification: $4^2 = 16$ and $10^3 + 2(10) + 3 = 1023 \equiv 16 \pmod{19}$ \checkmark}
\end{solbox}

\vspace{1em}

\subsection{Q6(b) --- 2 marks: Why ECC/RSA More Vulnerable to Quantum}

\begin{qbox}{2024 Q6(b) [2 marks]}
  Why are ECC and RSA more vulnerable than symmetric algorithms to quantum attacks?
\end{qbox}

\begin{solbox}{Solution}
\textbf{RSA and ECC are completely broken by Shor's Algorithm:}
\begin{itemize}[nosep]
  \item RSA security relies on the \textbf{integer factorization problem} (factoring $n=pq$ is hard classically — exponential time).
  \item ECC security relies on the \textbf{Elliptic Curve Discrete Logarithm Problem (ECDLP)} (computing $k$ from $kP$ is hard classically).
  \item \textbf{Shor's Algorithm} (quantum) solves \emph{both} integer factorization and discrete logarithm problems in \textbf{polynomial time} — completely breaking RSA and ECC.
\end{itemize}

\medskip
\textbf{Symmetric algorithms (AES) are only weakened, not broken:}
\begin{itemize}[nosep]
  \item Symmetric crypto has no mathematical structure (factorization/discrete log) for Shor's to exploit.
  \item \textbf{Grover's Algorithm} can search an $N$-element space in $O(\sqrt{N})$ time — providing only a \textbf{quadratic speedup}.
  \item Effect: AES-128 $\to$ $\sim$64-bit equivalent security. This is manageable by \textbf{doubling key size}.
  \item AES-256 remains \textbf{quantum-resistant}.
\end{itemize}
\end{solbox}

\vspace{1em}

%% ================================================================
\section{Lecture Example (Lc\#6 Slide 29--34) for Practice}
%% ================================================================

\begin{algobox}{Worked Example from Lecture: $y^2 \equiv x^3 + 2x + 2 \pmod{17}$}
\textbf{(3,6)+(3,6) --- Point Doubling on a different curve/prime:}

$x_1=3$, $y_1=6$, $a=2$, $p=17$
\[
  \lambda = \frac{3(9)+2}{2(6)} \pmod{17} = \frac{29}{12} \pmod{17}
\]
$29 \bmod 17 = 12$ and $12^{-1} \bmod 17 = 10$ (since $12 \times 10 = 120 \equiv 1 \pmod{17}$)
\[
  \lambda = 12 \times 10 = 120 \bmod 17 = 1
\]
\[
  x_3 = 1 - 6 = -5 \bmod 17 = 12 \qquad y_3 = 1(3-12)-6 = -15 \bmod 17 = 2
\]
\textbf{Result: $(3,6)+(3,6) = (12,2)$} (matches lecture slide 34)

\medskip
\textit{Note: The exam uses $b=3$, $p=19$ while the lecture uses $b=2$, $p=17$ --- same method, different numbers.}
\end{algobox}

\vspace{1em}

%% ================================================================
\section{Exam Strategy}
%% ================================================================

\begin{keybox}{What to Write in 5 Minutes}
\begin{tabular}{@{}p{3cm}p{10cm}@{}}
\toprule
\textbf{Part} & \textbf{What to write} \\
\midrule
Q6(a) [3 marks] & Write: Given $P=(3,6)$, $a=2$, $p=19$. Show Steps 1--3 with every intermediate mod. Box the answer $2P=(10,4)$. \\
Q6(b) [2 marks] & One sentence each for RSA (Shor breaks factorization), ECC (Shor breaks ECDLP), symmetric (Grover only halves, AES-256 safe). \\
\bottomrule
\end{tabular}

\medskip
\textbf{Watch out:}
\begin{itemize}[nosep]
  \item Use $x_3 = \lambda^2 - 2x_1$ for doubling (NOT $\lambda^2 - x_1 - x_2$, which is for $P \neq Q$ addition).
  \item Always convert negative results mod $p$ by adding $p$ repeatedly.
  \item For $a^{-1} \bmod p$: find $k$ such that $ak \equiv 1 \pmod{p}$ by trial or Extended Euclidean.
\end{itemize}
\end{keybox}

\end{document}
